\subsection{拓扑环}

定义 2. 拓扑环是一个集合 A，其上带有环结构与一个拓扑，满足下列公理：

\((\mathrm{AT}_{\mathbf{I}})\). 映射 \((x,y) \to x + y\)（\(\mathbf{A} \times \mathbf{A}\) 到 \(\mathbf{A}\) 中的映射）是连续的。

\((\mathrm{AT}_{\mathbf{II}})\). 映射 \(x \to -x\)（A 到 A 中的映射）是连续的。

\((\mathrm{AT}_{\mathrm{III}})\). 映射 \((x,y)\to xy\)（\(\mathbf{A}\times \mathbf{A}\) 到 \(\mathbf{A}\) 中的映射）是连续的。

前两条公理表明 \(\mathbf{A}\) 的拓扑与其加法群结构相容（§ 1，第 1 节）。

若在集合 A 上给定了环结构与一个拓扑，则当它们满足公理 \((\mathrm{AT}_{\mathbf{I}})\)、\((\mathrm{AT}_{\mathbf{II}})\) 与 \((\mathrm{AT}_{\mathbf{III}})\) 时，称它们是相容的。

例. 1) 在任一环 A 上，离散拓扑与环结构相容。拓扑为离散的拓扑环称为离散环。

* 2) 我们将在第四章看到，有理直线 Q（相应地，实直线 R）的拓扑与 Q（相应地，R）的环结构相容。 *

在拓扑环中，每个左位似 \(x \rightarrow ax\)（相应地，每个右位似 \(x \rightarrow xa\)）都是连续的（且当 a 是 A 的单位时是同胚）。

由于我们可以恒等地写出

\[
x y - x _ {0} y _ {0} = (x - x _ {0}) (y - y _ {0}) + (x - x _ {0}) y _ {0} + x _ {0} (y - y _ {0})
\]

公理 \((\mathrm{AT}_{\mathrm{III}})\){[}结合 \((\mathrm{AT}_{\mathrm{I}})\) 与 \((\mathrm{AT}_{\mathrm{II}})\){]}就等价于下列两条公理的合取：

\((\mathrm{AT}_{\mathrm{III}a})\). 对任意 \(x_0\in \mathbf{A}\)，映射 \(x\to x_0x\) 与 \(x\rightarrow xx_0\) 在点 \(x = 0\) 连续

\((\mathrm{AT}_{\mathbf{III}\ b})\). 映射 \((x,y)\rightarrow xy\)（\(\mathbf{A}\times \mathbf{A}\) 到 A 中的映射）在点 (o, o) 连续。

由此我们可以导出一组充要条件，使滤子 \(\mathfrak{B}\)（\(o\) 在一个环 \(A\) 中的邻域滤子）必须满足它们，才能在 \(A\) 上定义一个与其环结构相容的拓扑：\(\mathfrak{B}\) 必须满足 § 1 的公理 \((\mathrm{GA}_{\mathrm{I}})\) 与 \((\mathrm{GA}_{\mathrm{II}})\)，并且还满足下列两条公理：

\((\mathrm{AV}_{1})\). 对所有 \(x_0\in \mathbf{A}\) 与所有 \(\mathbf{V}\in \mathfrak{B}\)，存在 \(\mathbf{W}\in \mathfrak{B}\)，使得 \(x_0\mathbf{W}\subset \mathbf{V}\) 且 \(\mathbf{W}x_0\subset \mathbf{V}\)

\((\mathrm{AV}_{\mathrm{II}})\). 对所有 \(\mathbf{V} \in \mathfrak{B}\)，存在 \(\mathbf{W} \in \mathfrak{B}\)，使得 \(\mathbf{WW} \subset \mathbf{V}\)。

注. 在分析中我们相当经常遇到满足公理 \((\mathrm{AT}_{\mathrm{I}})\)、\((\mathrm{AT}_{\mathrm{II}})\) 与 \((\mathrm{AT}_{\mathrm{III}a})\)，但不满足 \((\mathrm{AT}_{\mathrm{III}b})\) 的环。* 紧群上的测度环就是一个例子，其中乘法是卷积，拓扑是淡拓扑。 *

例 3). 设 \(\mathfrak{B}\) 是环 A 上的一个滤子基，由双边理想组成。\(\mathfrak{B}\) 是某个拓扑（与 A 的加法群结构相容的拓扑）的 o 的基本邻域系，且由 \((\mathrm{AV}_{\mathrm{I}})\) 与 \((\mathrm{AV}_{\mathrm{II}})\) 立即得出，这个拓扑与 A 的环结构相容。

设 X 是拓扑空间，f 与 g 是 X 到拓扑环 A 中的两个映射。若 f 与 g 在一点 \(x_{0} \in X\) 连续，则 \(f + g\)、-f 与 fg 在该点连续。由此得出，X 到 A 中的连续映射组成 X 到 A 的一切映射所成之环 \(A^{X}\) 的一个子环。我们还看到，若 A 交换，则每个 n 个变元、系数在 A 中且定义在 \(A^{n}\) 上的多项式都在 \(A^{n}\) 上连续。再者，设 f 与 g 是一个集合 X（由滤子 F 所滤化）到豪斯多夫拓扑环 A 中的两个映射；若 \(\lim_{\mathfrak{F}} f\) 与 \(\lim_{\mathfrak{F}} g\) 存在，则 \(\lim_{\mathfrak{F}} (f + g)\)、\(\lim_{\mathfrak{F}} (-f)\) 与 \(\lim_{\mathfrak{F}} (fg)\) 也存在，并且我们有（第一章，§7，第 4 节，命题 9，推论 1，及 §8，第 1 节，命题 1）

\[
\lim _ {\mathfrak {F}} (f + g) = \lim _ {\mathfrak {F}} f + \lim _ {\mathfrak {F}} g,\tag{1}
\]

\[
\lim _ {\mathfrak {F}} (- f) = - \lim _ {\mathfrak {F}} f,\tag{2}
\]

\[
\lim _ {\mathfrak {F}} (f g) = (\lim _ {\mathfrak {F}} f) (\lim _ {\mathfrak {F}} g).\tag{3}
\]

